\subsection{Act cones with one total row excess and the sixteen-arrow SCC case}

\begin{theorem}[One total row excess in a finite categorical base]
\label{C:one-row-excess-act-cups}
Let $S$ be a finite strongly connected category.  In its minimal
Rees ideal, let $r_c$ be the number of rows at object $c$.
Suppose
\[
 \sum_{c\in\operatorname{Ob}S}(r_c-1)=1.
\]
For every nonempty finite right $S$-act $F$, all rational
products of degree-two classes on the act cone
$B\mathcal C(S,F)$ vanish.  Here the act cone adds one sink
object, with hom-set $F(c)$ from $c$ to that sink.
\end{theorem}
\begin{proof}
Write $R=\mathbb QS$, and let $\mathbf Q$ denote the constant
module, with one copy of $\mathbb Q$ at every object.  Choose
an idempotent in the minimal ideal and average its finite
maximal group $H$:
\[
 p=|H|^{-1}\sum_{h\in H}h\in R.
\]
The module $P=Rp$ is projective.  At object $c$ it has one
basis vector for each of the $r_c$ minimal rows: multiplying
an arrow by $p$ averages its free $H$-coordinate and fixes
its minimal-ideal column.  Sending each of these basis vectors
to $1$ gives a natural surjection $P\to\mathbf Q$.
Its kernel $V$ is the objectwise row augmentation module.
The assumed total excess means that $V$ is one-dimensional
at a unique object $a$ and zero at all other objects.
Endomorphisms of $a$ act by a scalar character $\chi$.
Since its image is a finite multiplicative subset of
$\mathbb Q$, its values lie in $\{0,1,-1\}$.

Every set $F(c)$ is nonempty: a nonempty stalk propagates
along arrows in both directions in the strongly connected
base.  Let $L$ be the kernel of the objectwise augmentation
$\mathbb QF\to\mathbf Q$ of right modules, and let $D L$
denote its objectwise rational dual.  The act-cone bar and
vector-space duality give
\[
 H^2(B\mathcal C(S,F);\mathbb Q)
 \cong\operatorname{Ext}^1_R(\mathbf Q,D L).
\]
Applying $\operatorname{Hom}_R(-,D L)$ to
$0\to V\to P\to\mathbf Q\to0$ gives the natural
surjection
\[
 \operatorname{Hom}_R(V,D L)\longrightarrow
 \operatorname{Ext}^1_R(\mathbf Q,D L).
\]
Thus every degree-two class has a representative given by
a right-module map $\ell:L\to D V$.

Choose a base state $s_c\in F(c)$ and put
\[
 q_c(s)=\ell_c(e_s-e_{s_c}).
\]
This is zero off $a$ and is a rational-valued function at $a$.
For an arrow $f:c\to d$, write $s f$ for the right action
$F(d)\to F(c)$.  Naturality says that
\[
 q_c(s f)=D V(f)q_d(s)+q_c(s_d f).
\]
Consequently the image sets $T(c)=q_c(F(c))$ form an actual
finite quotient right act.  Off $a$ they are singletons.
Every cross-arrow map between them is constant, because
$D V(f)=0$ for distinct source and target objects.  At $a$
the maps have the affine form $x\mapsto\chi(f)x+b_f$.
If $T(a)$ is a singleton, the representing functional and
the class are zero.  Otherwise a slope-one map of the finite
set is the identity, a slope-minus-one map is its common
reflection, and a slope-zero map is constant.

The functional $\ell$ factors through the augmentation
module of $T$.  Naturality of the preceding connecting
surjection proves that the original class is pulled back
from the act cone on $T$ under the quotient functor
$\mathcal C(S,F)\to\mathcal C(S,T)$.

Extend $T(a)$ to its convex interval $I(a)$ and put $I(c)$
equal to a point for $c\ne a$.  Extend the maps affinely.
They satisfy the category composition equations on these
intervals because they agree on the spanning finite sets;
the singleton cases are immediate.  This is a strict
contravariant diagram of contractible spaces.  Its quotient
diagram $I(c)/T(c)$ has a wedge of circles at $a$ and a point
elsewhere.  Every cross map is a pointed constant map.

Taking degree two on every circle commutes with every
endomorphism at $a$: identity, common reflection, and
constants.  It also commutes with all cross maps.  Hence it
is a strict self-map of this pointed diagram.  The pointed
homotopy colimit of $I/T$ is the cofiber of
$\operatorname{hocolim}T\to\operatorname{hocolim}I$, and
$\operatorname{hocolim}I\to BS$ is an equivalence by
objectwise contractibility.  This is the required act cone.
The reduced cellular complex of the quotient diagram is
concentrated in degree one, so the induced map of its bar
complex is multiplication by two in every degree.  It acts
by two on every positive rational cohomology group of the
act cone.  Multiplicativity therefore gives
$2xy=4xy$, forcing every positive-degree rational product
on this quotient act cone to vanish.

Pulling back kills the square of the original arbitrary
degree-two class.  Polarization over $\mathbb Q$ proves
the assertion for every pair of degree-two classes on the
original cone.  This proof allows a nontrivial minimal group:
its only use was the projective averaging idempotent $p$.
\end{proof}

\begin{theorem}[Three-object SCC exclusion at sixteen arrows]
\label{C:sixteen-three-SCC}
A three-object finite category with at most sixteen arrows
and a nontrivial strongly connected component cannot have
classifying space homotopy equivalent to $\mathbb {CP}^2$.
\end{theorem}
\begin{proof}
The smaller totals have already been excluded.  If one object
is a retract of another, delete it; at least its identity
and two opposite cross arrows disappear, leaving at most
thirteen arrows.  We may assume there is no such retract.
In any nontrivial strongly connected block, this implies
that every object identity lies outside the minimal ideal.
Indeed, an identity in that ideal makes its corner a group;
composing any two opposite cross arrows and adjusting by
the inverse of their corner product exhibits a retract.

First suppose all three objects are strongly connected.
If the total number of rows or the total number of columns
is three, the corresponding averaged minimal representable
is the constant module, which is projective.  The category
is rationally acyclic and cannot have the proposed type.
Otherwise both totals are at least four.  The minimal ideal
then has at least sixteen arrows, in addition to the three
identities outside it, exceeding the budget.

There remains a strong two-object block $S$, a singleton
monoid $M$, and cross arrows in one direction.  Write $s,m,t$
for their sizes.  Here $s\geq6$ and $t\geq2$.
If $t=2$, the two cross-hom sets are singletons, so the
category is a categorical join.  Its positive rational
products vanish.  We may take $t\geq3$ and $s+m\leq13$,
and use opposition to make $S$ the source block.

If both total row and total column counts of $S$ are at
least three, its minimal ideal has at least nine arrows;
the two identities give $s\geq11$.  Hence $s\leq12$ and
$m\leq2$.  Its minimal group is trivial and both total
counts are exactly three.  Its total row excess is one.
The trivial rational modules of every monoid of order at
most two are projective.  The actual orbit-collapse functor
at the $M$ vertex therefore gives a rational cohomology-ring
equivalence to a finite right-$S$ act cone.  Apply
Theorem~\ref{C:one-row-excess-act-cups}.  This includes the
twelve-arrow strong block with an extra arrow outside its
minimal band; no deletion or classification of that arrow
is needed.

In the remaining case the total row or column count of
$S$ is two, and $BS$ is rationally acyclic by the projective
constant-module argument.  If $m\leq4$, both vertex blocks
are rationally acyclic.  Collapsing them in the correspondence
homotopy pushout gives a rational ring equivalence to a
suspension, so all products vanish.

Suppose $m\geq5$.  Then $s\leq8$.  The minimal group of
$S$ is trivial, and its total row and column counts are
either $(2,2)$, $(2,3)$, or $(3,2)$.  For $(2,2)$ the
constant modules on both sides are projective.  Collapsing
$S$ by actual cross-act coinvariants therefore gives a
rational ring equivalence to a finite act cone over $M$,
where $m\leq7$.  Its rational homology would be bounded
and its second cohomology one-dimensional; this contradicts
Lemma~\ref{C:nine-bounded-rank-one}.

For $(2,3)$ or $(3,2)$, the six-element band and the two
identities exhaust $s=8$.  Choose an object $b$ with one
minimal row, and another object $a$ in the block.  Any
minimal idempotent $e\in\operatorname{End}(b)$ fixes every
incoming cross arrow at $b$: all such arrows belong to the
minimal band and their target row is unique.  There are
no incoming arrows from outside this source block.  The
only unit of the corner is its identity.  Minimal-band
arrows $u:b\to a$ and $v:a\to b$ can be chosen with
$vu=e$.  Lemma~\ref{C:one-sided-corner-deletion} deletes
$b$, leaving a smaller category already excluded by
Theorem~\ref{C:CP2-lower16}.  This proves the last case.
\end{proof}
